\documentclass[main.tex]{subfiles}
\begin{document}
\ifSubfilesClassLoaded{\begin{center}\small\itshape Part of the GnssProcessor Mathematical Reference \textbullet\ \today\end{center}}{}
\section{Orbit and Clock Product Evaluation}\label{sec:products}
% ============================================================
%  The deterministic satellite orbit/clock backbone that every
%  mode evaluates; the SSR overlay (sec:ssr) refines it.
% ============================================================

% doc-local notation (descriptive names to dodge package collisions)
\newcommand{\eciframe}{\bm{r}^{\sat}_{\mathrm{I}}}     % near-inertial derotated position
\newcommand{\rotz}[1]{\bm{R}_z(#1)}                    % rotation about ECEF z
\newcommand{\earthrot}{\omega_{\oplus}}               % Earth rotation rate (scalar)
\newcommand{\sidereal}{\kappa}                         % sidereal-day ratio
\newcommand{\trigfreq}{\Omega_{0}}                     % fundamental sidereal angular rate
\newcommand{\satclk}{\delta t^{\sat}}                 % satellite clock bias [s]
\newcommand{\eccanom}{E}                              % eccentric anomaly
\newcommand{\meananom}{M}                             % mean anomaly
\newcommand{\trueanom}{\nu}                           % true anomaly
\newcommand{\meanmotion}{n}                           % corrected mean motion
\newcommand{\jtwo}{J_2}                               % oblateness coefficient

Every processing mode in the engine (post-processed PPP, PPP-RTK, the
baseline modes, and the real-time streaming engine) shares one
deterministic satellite \emph{backbone}: a position
$\satpos$, velocity $\satvel$, and clock bias $\satclk$ for each
satellite at each query epoch.  This backbone is evaluated \emph{before}
any state-space refinement.  In post-processing it is the final word; in
the live engine it is the broadcast baseline that the SSR overlay
(\dcref{sec:ssr}{the SSR corrections document}) refines by a RAC orbit
delta and a scalar clock delta.  This document derives the backbone
itself, the two precise interpolation paths and the two broadcast
evaluation paths, plus the ephemeris-selection logic that feeds them.
It deliberately stops at the point where the SSR delta is added:
the orbit/clock \emph{overlay} math (RAC frame, IOD gate, freshness)
lives in \dcref{sec:ssr}{the SSR document} and is not re-derived here.
The Sagnac (Earth-rotation) range correction that turns $\satpos$ at
transmit time into a consistent line-of-sight belongs to
\dcref{sec:obsphys}{the observation-physics document}; the ERP /
polar-motion transforms that carry an SP3 sample between the terrestrial
and celestial frames belong to \dcref{sec:frames}{the reference-frames
document}.

The backbone splits cleanly into a \emph{precise} path and a
\emph{broadcast} path.  The precise path consumes tabulated SP3 orbit
samples and CLK clock samples (IGS Final / Rapid products) and
\emph{interpolates}; it is the post-processing default and is
centimetre-to-millimetre grade.  The broadcast path evaluates a
Keplerian (or, for GLONASS, a numerically integrated PZ-90) ephemeris
analytically from the navigation message; it is metre-grade on its own
and is the substrate the live SSR stream corrects.  On the default path
the two paths do not mix within a satellite at a given epoch: a precise
product, when loaded and spanning the query, wins; otherwise the engine
falls through to the broadcast evaluation
(\texttt{src/core/model/ModelGeometry.cpp},
\texttt{src/core/corrections/GnssCorrections.cpp}).  The resolution is
per quantity, however, not structurally coupled: the real-time RTB-4
gate \texttt{realtime.use\_precise\_orbits} splits the orbit and clock
sources independently (\texttt{GnssConfig::realtimeUsesLiveOrbitSource}
vs \texttt{realtimeUsesLiveOrbitClockSource},
\texttt{ProductParser::loadProducts(bool,\,bool)}), so an armed
``ultra-rapid orbits $+$ network clocks'' config evaluates one satellite
at one epoch with a precise interpolated orbit \emph{and} a
broadcast$+$SSR clock.  The flag defaults false, so the orbit gate then
returns exactly the clock gate, byte-identical, and the no-mixing
statement holds for every non-real-time mode and for the default
real-time path.

\subsection{Precise orbit: derotated Neville interpolation}\label{sec:prod-neville}

SP3 files tabulate the satellite position in the ECEF frame on a coarse
grid (typically $900$\,s for IGS Final, $300$\,s for Rapid).  The ECEF
trajectory is \emph{not} a smooth low-degree polynomial: it carries the
sidereal-frequency oscillation of the rotating frame relative to the
satellite's near-inertial orbit, which a degree-$10$ fit cannot capture
to better than a few millimetres near the SP3 epoch boundaries.  The
engine therefore \emph{derotates} each node into a common
near-inertial frame anchored at the query epoch before interpolating
(\texttt{src/core/model/ModelGeometry.cpp}, following the
RTKLIB / cssrlib convention \cite{takasu2009rtklib}).  For a node sampled
at $t_{\sat}$ and a query at $t_q$, with the Earth-rotation rate
$\earthrot = 7.2921151467\times10^{-5}$\,rad/s
(\texttt{Constants::EARTH\_ANGULAR\_VELOCITY}), the in-plane coordinates
are rotated by the angle the frame swept between the two epochs:
\begin{equation}
\eciframe(t_{\sat})
=
\rotz{\earthrot\,(t_{\sat}-t_q)}\;\satpos(t_{\sat}),
\qquad
\rotz{\alpha}=
\begin{bmatrix}
\cos\alpha & -\sin\alpha & 0\\
\sin\alpha & \phantom{-}\cos\alpha & 0\\
0 & 0 & 1
\end{bmatrix}.
\label{eq:prod-derotate}
\end{equation}
Concretely, with $\alpha_{\sat}=\earthrot\,(t_{\sat}-t_q)$ the derotated
in-plane components written to the interpolation arrays are
$\tilde{x}_{\sat}=\cos\alpha_{\sat}\,x_{\sat}-\sin\alpha_{\sat}\,y_{\sat}$
and
$\tilde{y}_{\sat}=\sin\alpha_{\sat}\,x_{\sat}+\cos\alpha_{\sat}\,y_{\sat}$,
with $z$ untouched (\texttt{ModelGeometry.cpp}).  Because the rotation
angle is exactly zero at $t_{\sat}=t_q$, the interpolated value at the
query is \emph{already} the ECEF position at $t_q$, so no post-rotation
is needed.

\paragraph{Bracket-centred window.}
The node window is the decisive design choice for position
\emph{continuity}.  The engine selects an $\mathrm{IGS\_POINTS}=11$-node
window centred on the \emph{bracket} $[t_a,t_{a+1}]$ that contains $t_q$,
not the $11$ nodes nearest to $t_q$
(\texttt{ModelGeometry.cpp}).  The distinction is subtle but
operationally important: a nearest-$N$ rule swaps a node at the
\emph{midpoint} between SP3 epochs, where the two candidate polynomials
disagree by a few millimetres, producing the well-known
synchronous-across-constellation $15$-minute sawtooth in PPP position.
A bracket-centred rule instead swaps the window exactly when $t_q$
crosses a node $t_{a+1}$; at that instant $t_{a+1}$ is shared by both
windows, so the Neville value is identical on either side and the
interpolated orbit is continuous across the swap.  A gap guard
re-picks a narrower $\mathrm{MIN\_POINTS}=9$-node window and rejects
the satellite if the residual span still exceeds the fit interval
(thresholds in \texttt{ModelGeometry.cpp}; tuning values cited, not
reproduced).  The ``fit interval'' is not read from the product header
but computed as the \emph{median} of the leading consecutive-node
spacings (\texttt{ModelGeometry.cpp}); this is deliberate, because a
multi-file SP3 merge can interleave a finer margin file's off-grid nodes
into a coarser arc (e.g. a $5$-minute ultra-rapid file concatenated with
a $15$-minute one, since the $1$\,s dedup window correctly keeps them as
non-duplicates, so the merged series is non-uniformly sampled and
\texttt{ProductParser} warns on the mixed-rate concat).  The median grid
gives a stable $\mathrm{MAX\_SPAN}$ guard on such a series where the
nominal grid would be ill-defined.

\paragraph{Neville value and derivative.}
Within the window the engine runs Neville's algorithm
simultaneously for the value and its analytic derivative
(\texttt{src/core/math/Interpolate.cpp},
\texttt{Interpolate::nevilleInterpDeriv}).  To avoid catastrophic
cancellation, the node abscissae are shifted to
$s_{\sat}=t_{\sat}-t_q$ (raw continuous GPS seconds are $\sim\!1.36\times10^8$,
so an un-shifted difference loses $\sim\!8$ digits); the polynomial value
and derivative are both invariant under this unit-Jacobian shift.  With
the tableau column $P_{i}$ for the value and $D_{i}$ for the derivative,
the coupled recurrence over degree $j$ (evaluating at the shifted query
$s_q=0$) is
\begin{align}
P_{i,j}
&=
\frac{(s_q-s_{i+j})\,P_{i,j-1}-(s_q-s_i)\,P_{i+1,j-1}}
     {s_i-s_{i+j}},
\label{eq:prod-neville-value}\\[4pt]
D_{i,j}
&=
\frac{(s_q-s_{i+j})\,D_{i,j-1}+P_{i,j-1}
       -(s_q-s_i)\,D_{i+1,j-1}-P_{i+1,j-1}}
     {s_i-s_{i+j}},
\label{eq:prod-neville-deriv}
\end{align}
seeded by $P_{i,0}=y_i$ and $D_{i,0}=0$; the converged scalars are
$P_{0,n-1}$ and $D_{0,n-1}$ (\texttt{Interpolate.cpp}).  The derivative
recurrence \eqref{eq:prod-neville-deriv} is exactly
$\partial/\partial s_q$ of the value recurrence
\eqref{eq:prod-neville-value}, so no second pass or finite difference is
required.

\paragraph{Inertial-to-ECEF velocity.}
The Neville derivative is taken on the \emph{derotated} channels, so it
is the velocity in the frame fixed at $t_q$, a snapshot inertial
frame.  Downstream consumers (Doppler line-of-sight rate, orbit-normal
attitude, the relativistic $\satpos\cdot\satvel$ term) require the ECEF
velocity, recovered by removing the frame rotation
$\satvel_{\mathrm{ECEF}}=\satvel_{\mathrm{I}}-\bm{\omega}_{\oplus}\times\satpos$.
For $\bm{\omega}_{\oplus}=(0,0,\earthrot)$ this is the boxed transform
applied in \texttt{ModelGeometry.cpp}:
\begin{empheq}[box=\fbox]{equation}
\satvel_{\mathrm{ECEF}}
=
\begin{bmatrix}
v_{\mathrm{I},x}+\earthrot\,\satpos_y\\[2pt]
v_{\mathrm{I},y}-\earthrot\,\satpos_x\\[2pt]
v_{\mathrm{I},z}
\end{bmatrix}.
\label{eq:prod-vel-ecef}
\end{empheq}
The relativistic clock term is invariant to this conversion, since
$\satpos\cdot(\bm{\omega}_{\oplus}\times\satpos)=0$; only the Doppler and
cross-product attitude paths see the difference
(\texttt{ModelGeometry.cpp}).

\subsection{Precise orbit: sidereal-trigonometric LSQ fit}\label{sec:prod-trig}

A config-gated alternative (\texttt{config.enableTrigSp3Interp()}) replaces
the Neville polynomial with an over-determined least-squares fit on a
\emph{sidereal-trigonometric basis}, run on the identical derotated
channels and time base
(\texttt{ModelGeometry.cpp};
\texttt{Interpolate::trigSvdInterp}).  The basis is a constant plus four
sine/cosine harmonic pairs of the fundamental sidereal angular rate,
nine terms in all:
\begin{equation}
\satpos_{\mathrm{c}}(\tau)
=
c_0
+
\sum_{f=1}^{4}\bigl[
  c_{2f-1}\sin(f\,\trigfreq\,\tau)
  + c_{2f}\cos(f\,\trigfreq\,\tau)
\bigr],
\qquad
\tau = t-t_0,
\label{eq:prod-trig-basis}
\end{equation}
where $t_0$ is the first node (the constant term absorbs the offset, the
sin/cos pairs the phase, so the fitted curve is invariant to this
choice).  The fundamental rate is the angular frequency of the
\emph{mean sidereal day}: with the sidereal-day ratio
$\sidereal=0.99726956634$ and a nominal solar day of $86400$\,s,
\begin{empheq}[box=\fbox]{equation}
\trigfreq
=
\frac{2\pi}{\sidereal\cdot 86400\ \mathrm{s}}\quad[\mathrm{rad/s}],
\qquad
\sidereal = 0.99726956634,
\label{eq:prod-trig-rate}
\end{empheq}
expressed \emph{per second} so the basis argument $\tau$ stays in the
caller's seconds and the analytic derivative falls out per second
(\texttt{Interpolate.cpp}; M-PAGES uses rad/day and divides the velocity
by $86400$, numerically identical).  The constant
$\sidereal=0.99726956634$ is the IERS mean sidereal-to-solar day ratio
and \emph{is} the physics of the model, so it is reproduced here with its
source line.

The fit is solved by Eigen's \texttt{JacobiSVD} with a relative
singular-value threshold of $10^{-15}$, which yields the minimum-norm
least-squares solution, the Numerical-Recipes \texttt{svbksb}
convention M-PAGES' \texttt{svdfit} uses \cite{press2007numerical}.  Per-node
weights derive from the SP3 \texttt{++} accuracy codes; because those are
a single value per satellite, the weight set is uniform and the fit is
effectively unweighted, matching M-PAGES.  The analytic value and
derivative at the query $\tau_q=t_q-t_0$ are
\begin{align}
\satpos_{\mathrm{c}}(\tau_q)
&=
c_0+\sum_{f=1}^{4}\bigl[
  c_{2f-1}\sin(f\trigfreq\tau_q)+c_{2f}\cos(f\trigfreq\tau_q)\bigr],
\label{eq:prod-trig-value}\\[4pt]
\dot{\satpos}_{\mathrm{c}}(\tau_q)
&=
\sum_{f=1}^{4}(f\trigfreq)\bigl[
  c_{2f-1}\cos(f\trigfreq\tau_q)-c_{2f}\sin(f\trigfreq\tau_q)\bigr],
\label{eq:prod-trig-deriv}
\end{align}
with the per-harmonic angular weight $w_f=f\,\trigfreq$
(\texttt{Interpolate.cpp}).  Crucially this is an \emph{over-determined
smoother}: the window is up to $11$ nodes against $9$ basis terms, so for
$N>9$ the fitted curve does \emph{not} pass through the nodes, a
deliberate divergence from the M-PAGES $9$-point exact fit, kept so the
trig branch reuses the existing derotated window verbatim.  The same
inertial-to-ECEF velocity conversion \eqref{eq:prod-vel-ecef} is then
applied to $\dot{\satpos}_{\mathrm{c}}$.

\subsection{Precise clock interpolation}\label{sec:prod-clock}

The CLK product tabulates the scalar satellite clock bias $\satclk$ on
its own grid (commonly $30$\,s).  Because clock noise is process-driven
(a near-Brownian random walk in bias, with occasional rate
discontinuities), the engine offers three interpolation modes selected by
\texttt{config.getSatelliteClockInterpolationMode()}
(\texttt{src/core/corrections/GnssCorrections.cpp}).  All three are
bracket-centred for the same continuity reason as the orbit window: the
sliding window must swap on a shared node, never at a between-epoch
midpoint.

\paragraph{(a) Linear (MMSE for Brownian motion).}
The default is a two-node linear interpolation across the bracket
$[t_0,t_1]$ containing $t_q$:
\begin{equation}
\satclk(t_q)
=
\satclk_0
+
\bigl(\satclk_1-\satclk_0\bigr)\,
\frac{t_q-t_0}{t_1-t_0},
\label{eq:prod-clock-linear}
\end{equation}
which is the minimum-mean-square-error estimator for a Wiener process
sampled at the two bracketing epochs.  Interpolation across a gap larger
than a guard threshold (in seconds) is rejected outright rather than
fabricated (\texttt{GnssCorrections.cpp}; gap value cited, not
reproduced).

\paragraph{(b) Rate-anchored weighted LSQ.}
The \texttt{LsqSatClock} mode fits a local line
$\satclk(t)\approx a+b\,(t-t_q)$ over an $\sim\!11$-node window by
weighted least squares, with the centring shift $s_k=t_k-t_q$ making the
fitted value at the query equal to the intercept $a$ directly.  The key
refinement is that the CLK product may also carry a clock \emph{rate}
$m_j$ at the coarse $5$-minute epochs; each rate is added as an
\emph{explicit observation of the slope}, $m_j=b+\mu_j$.  Since a value
equation $y_k=a+b\,s_k+\nu_k$ has $\partial/\partial a=1$, $\partial/\partial b=s_k$
while a rate equation $m_j=b+\mu_j$ has $\partial/\partial a=0$,
$\partial/\partial b=1$, the rate observations contribute to the slope
row only.  With value weights $w^{(y)}_k=1/\sigma_{y,k}^2$ and rate
weights $w^{(m)}_j=1/\sigma_{m,j}^2$, the $2\times2$ weighted normal
equations are
\begin{empheq}[box=\fbox]{equation}
\begin{bmatrix}
\displaystyle\sum_k w^{(y)}_k
&
\displaystyle\sum_k w^{(y)}_k s_k
\\[10pt]
\displaystyle\sum_k w^{(y)}_k s_k
&
\displaystyle\sum_k w^{(y)}_k s_k^{2}
\;+\;
\displaystyle\sum_j w^{(m)}_j
\end{bmatrix}
\!\!
\begin{bmatrix} a\\[6pt] b \end{bmatrix}
=
\begin{bmatrix}
\displaystyle\sum_k w^{(y)}_k y_k
\\[10pt]
\displaystyle\sum_k w^{(y)}_k s_k y_k
\;+\;
\displaystyle\sum_j w^{(m)}_j m_j
\end{bmatrix},
\label{eq:prod-clock-lsq}
\end{empheq}
solved in closed form for the intercept
\begin{equation}
a
=
\frac{\bigl(S_{wxx}+S_m\bigr)\,S_{wy}
       -S_{wx}\,\bigl(S_{wxy}+S_{my}\bigr)}
     {S_w\bigl(S_{wxx}+S_m\bigr)-S_{wx}^{2}},
\label{eq:prod-clock-lsq-soln}
\end{equation}
with the running sums named exactly as in
\texttt{GnssCorrections.cpp}
($S_w=\sum w^{(y)}_k$, $S_{wx}=\sum w^{(y)}_k s_k$,
$S_{wxx}=\sum w^{(y)}_k s_k^2$, $S_{wy}=\sum w^{(y)}_k y_k$,
$S_{wxy}=\sum w^{(y)}_k s_k y_k$, $S_m=\sum w^{(m)}_j$,
$S_{my}=\sum w^{(m)}_j m_j$), and the determinant in the denominator
guarding a degenerate (coincident-time) window with a linear fallback.
The estimate $\satclk(t_q)=a$ is returned.  Because the fit smooths
$\sim\!N$ value samples instead of touching two, the observation-noise
component of the estimate at the bracket midpoint drops by
$\sim\!\sqrt{N/2}$ relative to plain linear \eqref{eq:prod-clock-linear},
material in the seconds-to-minutes regime where the CLK observation noise
dominates the per-sample process noise.  The cost is a loss of
exact-node passage: the smoother introduces millimetre-level
discontinuities at the clock epochs when the window slides.

\paragraph{(c) Barycentric Lagrange.}
For high-rate ($1$\,Hz) products a barycentric-form Lagrange interpolation
over a $\mathrm{LAGRANGE\_POINTS}=9$-node window is available
(\texttt{GnssCorrections.cpp};
\texttt{Interpolate::barycentricWeights} / \texttt{barycentricEval}).  The
barycentric form
\begin{equation}
\satclk(t_q)
=
\Biggl.
\sum_{i=1}^{n}\frac{w_i}{t_q-t_i}\,\satclk_i
\;\Biggm/\;
\sum_{i=1}^{n}\frac{w_i}{t_q-t_i},
\qquad
w_i=\prod_{k\neq i}\frac{1}{t_i-t_k},
\label{eq:prod-clock-bary}
\end{equation}
is the numerically stable evaluation of the degree-$(n-1)$ Lagrange
polynomial \cite{berrut2004barycentric}; an extrapolation guard falls back
to two-node linear when $t_q$ lies outside the node span.  This mode is
\emph{not} recommended for standard $30$\,s IGS Final clocks, where the
polynomial oscillates between the sparse samples, and is retained only
for genuinely high-rate input.

\subsection{Broadcast Keplerian evaluation}\label{sec:prod-kepler}

When no precise product spans the query, GPS, Galileo, BeiDou and QZSS
are evaluated from their Keplerian broadcast ephemeris
(\texttt{src/core/products/LiveEphemerisStore.cpp},
\texttt{keplerPosition}).  This is the substrate the live SSR stream
corrects.  Each system carries its own gravitational parameter $\mu$ and
Earth-rotation rate $\earthrot$ (\texttt{LiveEphemerisStore.cpp}):
GPS / QZSS use $\mu=3.986005\times10^{14}\,\mathrm{m^3/s^2}$,
$\earthrot=7.2921151467\times10^{-5}\,\mathrm{rad/s}$;
Galileo / BeiDou use $\mu=3.986004418\times10^{14}\,\mathrm{m^3/s^2}$,
with BeiDou taking the slightly different
$\earthrot=7.292115\times10^{-5}\,\mathrm{rad/s}$.  BeiDou is evaluated
entirely on its native BeiDou Time (BDT) scale, where BDT week $0$ aligns
with GPS week $1356$ and $\mathrm{BDT}=\mathrm{GPS}-14$\,s, with both
the query epoch and the stored $t_{oe}$ kept on BDT so the node-rotation
term below uses the BDT seconds-of-week directly
(\texttt{LiveEphemerisStore.cpp}).

With the semi-major axis $a=(\sqrt{a})^2$ and the time-from-ephemeris
$t_k=t-t_{oe}$ reduced into $[-302400,+302400]$\,s (week rollover), the
corrected mean motion adds the broadcast $\delta n$ to the Keplerian
value, and the mean anomaly propagates linearly:
\begin{equation}
\meanmotion
=
\sqrt{\frac{\mu}{a^{3}}}+\delta n,
\qquad
\meananom
=
M_0+\meanmotion\,t_k.
\label{eq:prod-kepler-motion}
\end{equation}
Kepler's equation $\meananom=\eccanom-e\sin\eccanom$ is solved by Newton
iteration (up to $15$ steps, convergence at
$|\Delta\eccanom|<10^{-13}$):
\begin{equation}
\eccanom_{k+1}
=
\eccanom_k
+
\frac{\meananom-\bigl(\eccanom_k-e\sin\eccanom_k\bigr)}
     {1-e\cos\eccanom_k}.
\label{eq:prod-kepler-newton}
\end{equation}
The true anomaly and the argument of latitude follow, then the broadcast
second-harmonic corrections refine the argument of latitude $u$, the
orbital radius $r$, and the inclination $i$:
\begin{align}
\trueanom
&=
\operatorname{atan2}\!\bigl(\sqrt{1-e^2}\,\sin\eccanom,\;\cos\eccanom-e\bigr),
\qquad
\phi=\trueanom+\omega,
\label{eq:prod-kepler-anom}\\[4pt]
u &= \phi + c_{us}\sin 2\phi + c_{uc}\cos 2\phi,
\nonumber\\
r &= a\bigl(1-e\cos\eccanom\bigr) + c_{rs}\sin 2\phi + c_{rc}\cos 2\phi,
\label{eq:prod-kepler-harmonic}\\
i &= i_0 + \dot{\imath}\,t_k + c_{is}\sin 2\phi + c_{ic}\cos 2\phi.
\nonumber
\end{align}
The corrected longitude of the ascending node folds in the node rate
\emph{relative to} the rotating Earth and the fixed node offset at the
ephemeris epoch:
\begin{empheq}[box=\fbox]{equation}
\Omega
=
\Omega_0
+
\bigl(\dot{\Omega}-\earthrot\bigr)\,t_k
-
\earthrot\,t_{oe},
\label{eq:prod-kepler-node}
\end{empheq}
where the final $-\earthrot\,t_{oe}$ uses the seconds-of-week of $t_{oe}$
on the native scale (\texttt{LiveEphemerisStore.cpp}).  With the in-plane
position $(x_o,y_o)=(r\cos u,\,r\sin u)$ the ECEF position is the standard
rotation through $\Omega$ and $i$:
\begin{equation}
\satpos
=
\begin{bmatrix}
x_o\cos\Omega - y_o\cos i\,\sin\Omega\\
x_o\sin\Omega + y_o\cos i\,\cos\Omega\\
y_o\sin i
\end{bmatrix}.
\label{eq:prod-kepler-ecef}
\end{equation}
The BeiDou GEO additional fixed rotation is not applied (GEO satellites
are rarely used for PPP-RTK).  Velocity is obtained by central finite
difference of \eqref{eq:prod-kepler-ecef} over $\pm0.5$\,s, $\satvel
\approx[\satpos(t+\tfrac12)-\satpos(t-\tfrac12)]/(2\,\Delta t)$ with
$\Delta t=0.5$\,s, reusing the same
analytic position evaluator (\texttt{LiveEphemerisStore.cpp}).

\subsection{Broadcast GLONASS: PZ-90 RK4 integration}\label{sec:prod-glonass}

GLONASS does not broadcast Keplerian elements; it broadcasts an ECEF
state vector (position, velocity, and a luni-solar acceleration) valid at
$t_b$, and the ICD prescribes numerical integration of the PZ-90
equations of motion to the query epoch \cite{glonassicd2008}.  The engine
marches a fixed-step fourth-order Runge--Kutta integrator
(\texttt{src/core/products/LiveEphemerisStore.cpp},
\texttt{glonassState}).  The per-axis acceleration combines the central
gravitational term, the $\jtwo$ oblateness perturbation, the centrifugal
and Coriolis terms of the rotating frame, and the broadcast luni-solar
acceleration held constant over the integration.  Writing
$r=\lVert\satpos\rVert$ and the oblateness coefficient
\begin{equation}
\kappa_{\jtwo}
=
\tfrac{3}{2}\,\jtwo\,\frac{\mu\,a_e^{2}}{r^{5}},
\qquad
\jtwo = 1.0826257\times10^{-3},
\label{eq:prod-glonass-j2coef}
\end{equation}
(with $\mu=3.9860044\times10^{14}\,\mathrm{m^3/s^2}$,
$a_e=6378136\,\mathrm{m}$, $\earthrot=7.292115\times10^{-5}\,\mathrm{rad/s}$,
all PZ-90 constants at \texttt{LiveEphemerisStore.cpp}), the accelerations are
\begin{empheq}[box=\fbox]{align}
\ddot{x}
&=
-\frac{\mu}{r^{3}}x
+\kappa_{\jtwo}\,x\bigl(1-5\,z^2/r^2\bigr)
+\earthrot^{2}x+2\earthrot\dot{y}+a_{x},
\nonumber\\
\ddot{y}
&=
-\frac{\mu}{r^{3}}y
+\kappa_{\jtwo}\,y\bigl(1-5\,z^2/r^2\bigr)
+\earthrot^{2}y-2\earthrot\dot{x}+a_{y},
\label{eq:prod-glonass-accel}\\
\ddot{z}
&=
-\frac{\mu}{r^{3}}z
+\kappa_{\jtwo}\,z\bigl(3-5\,z^2/r^2\bigr)
+a_{z},
\nonumber
\end{empheq}
where $(a_x,a_y,a_z)$ is the broadcast luni-solar acceleration.  Note the
distinct numerator on the $z$-axis oblateness term ($3-5z^2/r^2$ versus
$1-5z^2/r^2$ in-plane), and that the centrifugal $\earthrot^2$ and
Coriolis $2\earthrot$ terms appear only on the in-plane axes, as required
in the PZ-90 rotating frame.  The integrator steps with a fixed maximum
$h=30$\,s, signed toward the query, clamping the final partial step so
the march lands exactly on $t_q$; each step is the classical RK4 update
$y\!\leftarrow\! y+\tfrac{h}{6}(k_1+2k_2+2k_3+k_4)$
(\texttt{LiveEphemerisStore.cpp}).

\subsection{Broadcast satellite clock}\label{sec:prod-bclock}

The broadcast satellite clock for the Keplerian systems is the standard
second-order polynomial in the time-from-clock-epoch $t-t_{oc}$
(on the native scale, with BeiDou shifted to BDT exactly as the orbit
node term is):
\begin{empheq}[box=\fbox]{equation}
\satclk
=
a_{f0}+a_{f1}\,(t-t_{oc})+a_{f2}\,(t-t_{oc})^{2}.
\label{eq:prod-bclock-poly}
\end{empheq}
GLONASS instead uses its two-term broadcast model with the $1020$-message
$\tau_n$ (DF123) and $\gamma_n$ (DF120) about $t_b=t_{oe}$:
\begin{equation}
\satclk
=
-\tau_n+\gamma_n\,(t-t_b)
\label{eq:prod-bclock-glonass}
\end{equation}
(\texttt{LiveEphemerisStore.cpp}).  This path is exercised only where
GLONASS is excluded from \texttt{active\_systems}, i.e.\ the shipped
real-time replay configs, which is the context this store serves; the
default \texttt{active\_systems} and the post-processing configs both
enable GLONASS, so the path is not universally dead code.

\paragraph{The precise-clock-baseline convention.}
A design decision worth stating explicitly: the broadcast clock
\eqref{eq:prod-bclock-poly} \emph{deliberately omits} both the periodic
relativistic eccentricity term and the group delay (TGD/BGD).  This is
not an oversight: it is what makes the broadcast clock a clean drop-in
for the SSR scalar clock delta.  The IGS-RTS clock correction is
generated against precisely this bare polynomial, so the engine must
\emph{not} pre-apply the periodic relativity or TGD before adding the SSR
delta, or the correction would double-count.  The periodic relativistic
term
$-2\,\satpos\cdot\satvel/\sLight^2$ is added later, in the
observation-physics path; see
\dcref{sec:obsphys}{the observation-physics document} for its
derivation.  This convention keeps the orbit, the bare clock, and the SSR
overlay mutually consistent on one ephemeris baseline.

\subsection{Ephemeris selection and IOD bookkeeping}\label{sec:prod-selection}

The broadcast paths above presuppose that the \emph{right} ephemeris
record has been selected.  The live store holds, per satellite, a
\texttt{current} and an optional \texttt{previous} record, and several
interlocking rules govern which one feeds the orbit/clock evaluation
(\texttt{src/core/products/LiveEphemerisStore.cpp}).

\paragraph{Freshest-toe within a source.}
Within a single source, the record with the newer reference epoch wins:
an incoming record whose $t_{oe}$ is older than the held one (a stale
re-broadcast) is rejected, while a sub-second jitter or equal-epoch
refresh overwrites in place (\texttt{LiveEphemerisStore.cpp}).  Selection
between \texttt{current} and \texttt{previous} at query time is by
\emph{closest} $t_{oe}$ on a common GPS-continuous scale, where BeiDou's BDT
$t_{oe}$ is lifted to GPS seconds first
(\texttt{toeGpsContinuous}), or every BeiDou record would read $\sim\!1356$
weeks stale and be wrongly skipped.

\paragraph{IGS-source authority over receiver nav.}
When the engine ingests both an IGS broadcast stream and the receiver's
own decoded navigation, the IGS record is \emph{authoritative}.  The
IGS-RTS SSR delta is computed against IGS broadcast ephemeris, so the
rover must pair the delta with the \emph{same} IGS broadcast; even at an
identical IOD, the receiver's locally assembled orbit carries slightly
different fit parameters (a live Galileo I/NAV assembly was observed to
diverge by $100$\,m$+$), and applying IGS-SSR on top of a receiver
ephemeris detonates the single-difference frame.  A receiver record
therefore never overwrites a current IGS record at an equal-or-older
$t_{oe}$; it is accepted only when it is strictly newer than the IGS
record by more than one constellation refresh ($1800$\,s, the GLONASS
nominal step), the signature of a \emph{stalled} IGS feed, so a dead
IGS stream degrades gracefully to receiver nav rather than stranding the
satellite (\texttt{LiveEphemerisStore.cpp}).

\paragraph{IODE-change demote and SSR re-referencing.}
When the IOD changes, the \texttt{current} record is demoted to
\texttt{previous} (rather than discarded), keeping the just-superseded
orbit valid across the cutover instant
(\texttt{LiveEphemerisStore.cpp}).  This matters because, across a
cutover, the receiver can hold a version one step ahead of the lagging
SSR stream.  When an SSR delta names a specific IOD,
\texttt{pickForSsr} re-references: it returns whichever held record
(\texttt{current} or \texttt{previous}) matches the SSR's IOD, so the
orbit and clock are evaluated on the exact baseline the correction was
computed for, instead of dropping the satellite to broadcast grade on a
\texttt{current}-vs-SSR mismatch.  The same current/previous retention
underwrites a second consumer built on the broadcast clock polynomial
\eqref{eq:prod-bclock-poly}:
\texttt{LiveEphemerisStore::ephCutoverShiftMetres} evaluates
$\sLight\,(\satclk_{\mathrm{new}}-\satclk_{\mathrm{old}})$ across an
ephemeris version roll (\texttt{computeSatelliteClock} run once on the
freshest pick and once pinned to the old IODE), and returns nothing
when the old version is no longer retained.  It is the product-side half
of the network producer's \texttt{network.apriori\_source = brdc\_rts}
eph-cutover re-level; that re-level is derived in
\dcref{sec:net}{the Network mode document} and is not repeated here.

\paragraph{The GNSS-IOD matching convention.}
The SSR ``GNSS IOD'' field (IGS-SSR / RTCM \texttt{4076}, IDF012) is an
$8$-bit quantity whose per-constellation \emph{definition} differs from
the raw broadcast IODE stored on the ephemeris.  \texttt{ssrMatchIode}
normalises the held record's IODE into the SSR convention in exactly one
place (\texttt{LiveEphemerisStore.cpp}):
\begin{equation}
\mathrm{IOD}_{\mathrm{SSR}}(\sat)
=
\begin{cases}
\mathrm{IODE} & \text{GPS, QZSS (8-bit LNAV, passthrough),}\\[2pt]
\mathrm{IODnav}\ \&\ \mathtt{0xFF}
              & \text{Galileo (low 8 bits of the 10-bit I/NAV IODnav),}\\[2pt]
\bigl\lfloor t_{oe}^{\,\mathrm{sow}}/720\bigr\rfloor \bmod 240
              & \text{BeiDou (toe-derived, per IGS-SSR v1.00),}
\end{cases}
\label{eq:prod-ssr-iode-map}
\end{equation}
with $t_{oe}^{\,\mathrm{sow}}$ the BDT seconds-of-week.  The BeiDou
broadcast AODE that the decoder stores in the IODE field is a
\emph{different} quantity and would never match the SSR field, so it is
deliberately ignored in favour of the toe-derived value.  The raw IODE is
left intact for the demote bookkeeping above; only the value handed to
the SSR gate is normalised here.  How that normalised IOD then \emph{gates}
the application of the SSR orbit/clock delta (accept the delta only on
an exact IOD match, else fall back to broadcast grade) is the subject
of \dcref{sec:ssr-iod}{the SSR IOD-gate section}; this document supplies
only the matched quantity it consumes.

\ifSubfilesClassLoaded{\printbibliography}{}
\end{document}
